\scanpage{48}
% Pencil archival foliation: 47.
Sous la forme précédente, on voit que
l'on a mieux que \deleted{des} \unsure{invariants|invariant}
dans $\mathrm{M}^{+}_i(K)$. En fait, on a des
\[
\underline{\gamma}^{\,i}(X)\in
\operatorname{Ob}\mathcal{M}^{+}_{i\ \mathrm{bipur}}(K),
\]
objets semi-simples;
définis de la façon suivante ($K$ parfait,
ce qu'on peut supposer)
\inserted{$\gamma^i(X)\unsure{\subset|\in}
\varinjlim_{U\ \mathrm{ouvert\ dense}}H^i(U,\mathbb{Z}(0))$}
est le plus grand sous-motif \uline{bipur} de
type $i$.
\deleted{[Ainsi le foncteur $\unsure{\gamma^i(X)}$}
\inserted{C'est un foncteur en $X$}
[\unsure{via|si on} \unsure{l'anneau} des fonctions rat. sur $X$,
comme sp. sur $K$], de façon évidente
\margin{\diagram{\begin{tikzcd}X'\arrow[d]\\X\end{tikzcd}}}
[NB si $X'\unsure{-|\to}X$ n'est pas dominant,
$X$ \unsure{irrd|irréd}, alors
$\gamma^i(X)\unsure{-|\to}\gamma^i(X')$ est
nul]. Pour un hom \uline{dominant},
$\gamma^i(X)\to\gamma^i(X')$ est \uline{injectif}. En effet,
on peut supposer (par des modèles convenables)
que $X'\unsure{-|\to}X$ est propre, \inserted{donc surjectif}
$X$ et $X'$ lisses,
alors $H^i(X)\unsure{-|\to}H^i(X')$ est \uline{injectif} comme
il est bien connu ----

\uline{Question} \deleted{\unsure{La relation}}
Propriétés de
``fidélité'' \deleted{de} des foncteurs précédents ??
On aimerait
que tout \deleted{hom} système d'homs
$\underline{\unsure{\gamma|\beta}}_i(\unsure{R|K})\to
\underline{\gamma}_i(R')$
